% !PRINTED VII
\weberpage{VII}

\noindent nach Möglichkeit die wichtigsten Quellen und literarischen Hülfsmittel an geeigneter Stelle zu erwähnen.

Es ist mir eine angenehme Pflicht, so manchem Freunde und Collegen, der an dem Fortschreiten meiner Arbeit regen und thatkräftigen Antheil genommen hat, hier meinen Dank auszusprechen. Zuerst gilt dieser Dank meinem Freunde \gesperrt{Dedekind} für seine treue Hülfe bei der Correctur, und wenn er auch auf den Plan und die Ausführung meines Werkes keinen Einfluss ausgeübt hat, so möchte ich doch nicht unerwähnt lassen, dass ich schon vor vielen Jahren durch ein Heft einer Vorlesung, die er im Winter 1857/58 in Göttingen über höhere Algebra, insbesondere über die Theorie von \gesperrt{Galois} gehalten hat, ein noch lebhafteres Interesse für diese Theorie gewonnen habe, die vordem auf unseren Hochschulen in solcher Vollständigkeit wohl noch nicht vorgetragen war.

Auch der mannigfachen Anregung und Belehrung habe ich hier zu gedenken, die ich meinem Freund und Collegen F.~\gesperrt{Klein} verdanke, der das Fortschreiten der Arbeit mit regstem Interesse begleitet hat und dessen sachkundiger, stets bereitwilligst gegebener Rath in manchen Theilen des Buches von grossem Einfluss gewesen ist.

Ich kann hier nicht alle Fachgenossen und Freunde namhaft machen, die mich durch ihren Rath unterstützt haben, wie der Leser an den betreffenden Stellen finden wird. Aber der Herren E.~\gesperrt{Hess} in Marburg, Fr.~\gesperrt{Meyer} in Clausthal, R.~\gesperrt{Fricke} in Braunschweig, die durch kundige und sorgfältige Ausführung der mühevollen Correctur der Druckbogen Genauigkeit und Richtigkeit des Textes gefördert haben, muss ich hier noch gedenken.

Endlich gilt mein Dank der Verlagsbuchhandlung, die durch bereitwilliges Eingehen auf meine Wünsche, durch Sorgfalt in Druck und Ausstattung wesentlich zum Gelingen des Ganzen beigetragen hat.

\bigskip

\gesperrt{Göttingen}, im November 1894.

\medskip

\begin{flushright}
\textbf{\gesperrt{Der Verfasser}.}
\end{flushright}
